\begin{table*}[!t]
\centering
\caption{Module~4 constants. All rates except the PkbA arm (4.10--4.14) carry the uniform time-scale factor $s_4=20$ (Sec.~\ref{sec:m4}).}
\label{tab:m4const}
\footnotesize
\renewcommand{\arraystretch}{1.15}
\begin{tabularx}{\textwidth}{@{}l r l c >{\raggedright\arraybackslash}X@{}}
\toprule
Symbol & Value & Unit & Basis & Meaning, justification, source\\
\midrule
$s_4$ & \num{20} &  & C & Uniform clock multiplier on all Module~4 rates. Leaves every steady state and loop gain unchanged; it makes PIP3 formation complete within seconds of a 40\,s stimulus. \\
$[\mathrm{PI3K}]_\mathrm{tot}$ & \num{200000} & $\mathrm{molec/voxel}$ & E & Pool large enough that catalyst depletion is negligible. \\
$[\mathrm{PIP2}]_\mathrm{tot}$ & \num{769231} & $\mathrm{molec/voxel}$ & E & Lipid pool ($10^7$ per cell), conserved: $\mathrm{PIP2}+\mathrm{PIP3}=$ const. \\
$[\mathrm{PTEN}]_\mathrm{tot}$ & \num{10000} & $\mathrm{molec/voxel}$ & E & PTEN pool, conserved between cytosol and membrane. \\
$[\mathrm{PKBA}]_\mathrm{tot}$ & \num{100000} & $\mathrm{molec/voxel}$ & C & Pool of the PH-domain kinase PkbA; enlarged so that its feedback can act at low activated fractions. \\
$k_{3K,\mathrm{on}}$ & \num{0.4} & $\mu\mathrm{M}^{-1}\mathrm{s}^{-1}$ & C & \mkt{M4-ras-pi3k}{RasG-GTP recruits PI3K to the membrane}~\cite{sasaki2004,fukushima2019}. \\
$k_{3K,\mathrm{off}}$ & \num{10} & $\mathrm{s^{-1}}$ & C & PI3K membrane release; \mkt{M4-pip3-trans}{PI3K activity is transient after a stimulus}~\cite{huang2003}. \\
$k_\mathrm{syn}$ & \num{0.0955} & $\mu\mathrm{M}^{-1}\mathrm{s}^{-1}$ & C & PIP3 synthesis per membrane PI3K; sets the PIP3 count (shot noise), not its contrast. \\
$k_{PT,\mathrm{on}},\ k_{PT,\mathrm{off}}$ & \num{0.6},\ \num{4} & $\mu\mathrm{M}^{-1}\mathrm{s}^{-1},\ \mathrm{s^{-1}}$ & C & PTEN docking on PIP2 and release; \mkt{M4-recip}{PTEN sits on the membrane at the back of the cell}~\cite{funamoto2002,iijima2002}. \\
$k_\mathrm{ev}$ & \num{800} & $\mathrm{s^{-1}}$ & C & Saturated PIP3-dependent eviction rate of PTEN. \\
$K_\mathrm{ev},\ n_\mathrm{ev}$ & \num{1.31e5},\ 2 & $\mathrm{molec/voxel},\ -$ & C & Half-eviction level ($17\,\%$ of the lipid pool) and steepness of the eviction gate (Sec.~\ref{sec:hill}). \\
$k_{PT}$ & \num{1.39} & $\mu\mathrm{M}^{-1}\mathrm{s}^{-1}$ & C & \mkt{M4-pten-dephos}{PTEN dephosphorylation of PIP3 to PIP2}~\cite{iijima2002,huang2003}. \\
$k_\mathrm{SH}$ & \num{1.28} & $\mathrm{s^{-1}}$ & C & PTEN-independent 5-phosphatase removal of PIP3; $k_\mathrm{SH}=k_\mathrm{ship}[\mathrm{SHIP}]$ with a constant catalyst pool of 77 molecules/voxel. \\
$k_{K,\mathrm{on}}$ & \num{1.32e-3} & $\mu\mathrm{M}^{-1}\mathrm{s}^{-1}$ & D & PH-domain recruitment of PkbA; solved so that the resting rate is $0.002\,\mathrm{s^{-1}}$ at $7298$ PIP3/voxel. \\
$k_{K,\mathrm{off}}$ & \num{0.05} & $\mathrm{s^{-1}}$ & E & PkbA membrane release. \\
$k_{K,\mathrm{act}}$ & \num{0.507} & $\mu\mathrm{M}^{-1}\mathrm{s}^{-1}$ & D & TORC2-dependent activation; solved so that the resting rate is $8.5\times10^{-3}\,\mathrm{s^{-1}}$ at 81 TORC2$^*$/voxel. \\
$k_{K,\mathrm{deact}}$ & \num{0.02} & $\mathrm{s^{-1}}$ & E & PkbA$^*$ decay ($\tau=50$\,s); sets the refractory duration. \\
$k_{K,\mathrm{fb}}$ & \num{160} & $\mu\mathrm{M}^{-1}\mathrm{s}^{-1}$ & C & Strength of the PKB feedback on PI3K, chosen so that PIP3 returns to baseline after a pulse. \\
$D_{\mathrm{PI3K}_c},D_{\mathrm{PTEN}_c},D_{\mathrm{PKBA}_c}$ & \num{10} & $\mu\mathrm{m^2/s}$ & E & Cytosolic value; the PTEN shuttle crosses the cell in $\approx\!2.5$\,s. \\
$D_{\mathrm{PI3K}_m},D_{\mathrm{PTEN}_m},D_{\mathrm{PIP2/3}},D_{\mathrm{PKBA}_m}$ & \num{0.1} & $\mu\mathrm{m^2/s}$ & E & Membrane-bound species; the active PkbA is kept local so that feedback ends the patch where it forms. \\
\bottomrule
\end{tabularx}
\end{table*}
